\documentclass[a4paper]{article}
\usepackage{graphicx}
\usepackage{amsmath,amssymb}
\usepackage{minted}
\usepackage{multirow}
\usepackage{float}
\usepackage{algorithm}
\usepackage{algpseudocode}
\usepackage{todonotes}
\usepackage{forest}
\usepackage{xstring}
\usepackage{xcolor}
\usepackage[most]{tcolorbox}
\usepackage{xparse}
\usepackage{natbib}
\usetikzlibrary{graphs,graphdrawing}
\usegdlibrary{layered}
% ND Mind Mirror default input files.
\input{/home/donkarlo/Dropbox/repo/nd_language_ai_project/src/nd_language_ai/formal/computer/programming/tex/latex/code_clips/external_graphviz.tex}
\input{/home/donkarlo/Dropbox/repo/nd_language_ai_project/src/nd_language_ai/formal/computer/programming/tex/latex/part/control_sequence/kind/semtag/semtag.tex}
\input{/home/donkarlo/Dropbox/repo/nd_language_ai_project/src/nd_language_ai/formal/computer/programming/tex/latex/code_clips/seehere.tex}

\title{Trend Detection in Time Series}
\date{}
\author{Mohammad Rahmani}
\begin{document}
\maketitle

\section{Definition}
Trend detection identifies persistent changes in the underlying level of a time series. It differs from forecasting: trend analysis describes the past; forecasting predicts future values.

An exploratory additive decomposition is
\[
 y_t = T_t + S_t + R_t,
\]
where $T_t$ denotes trend, $S_t$ denotes seasonality, and $R_t$ is the remainder. An additive decomposition is not always appropriate: positive retail data may instead exhibit multiplicative seasonality.

\section{Common Methods}
\begin{itemize}
 \item Moving Average: interpretable smoothing of short-run fluctuations.
 \item Exponential Moving Average: online smoothing emphasizing recent data.
 \item Linear Regression: interpretable linear trend and slope.
 \item Mann--Kendall Test: rank-based significance test for monotonic trend.
 \item Theil--Sen Estimator: robust estimate of linear trend slope.
 \item STL Decomposition: separate seasonal, trend, and remainder components.
 \item Change-Point Detection: discover structural breaks in level or slope.
 \item LOESS/LOWESS: flexible local estimation of nonlinear trends.
\end{itemize}
Each method has a standalone LaTeX document in its own subfolder.

\section{E-Commerce Analysis Workflow}
\begin{enumerate}
 \item Define a daily or weekly measure: revenue, units sold, or conversion rate.
 \item Normalize timestamps and distinguish missing values from real zeros.
 \item Plot raw data; annotate holidays, promotions, and stockouts.
 \item Use a moving average or EMA for initial exploratory smoothing.
 \item Account for day-of-week and seasonal variation using STL if appropriate.
 \item Quantify a roughly linear trend with regression or Theil--Sen.
 \item Test for monotonic trend using Mann--Kendall only after checking dependence.
 \item Investigate regime changes before reporting a single global trend.
 \item For forecasting, fit every transformation using only the training history.
\end{enumerate}

\section{Important Distinctions}
Smoothing does not establish statistical significance. A Mann--Kendall test checks for evidence of a monotonic direction, whereas Theil--Sen estimates its magnitude. A change-point method identifies structural transitions rather than one continuous slope. Two-sided smoothers may use future observations and cause data leakage if used as prediction-time features.

\bibliographystyle{plainnat}
\bibliography{/home/donkarlo/Dropbox/repo/nd_shared_research/bibliography/bibliography.bib}
\end{document}
